\chapter{Programmablaufprotokolle}

\section{Steuerung}

Die Spielfigur wird durch das anklicken der gelben Pfeile an den Rändern des Spielfensters bzw. der Türen durch das Schulhaus bewegt. Wird die Spielfigur selbst angeklickt, öffnet sich deren Inventar
und bereits gesammelte Items können eingesehen werden. Das Inventar kann mit einem Klick auf sein rotes \glqq{}X\grqq{} wieder geschlossen werden.

Wird die Taste \texttt{P} gedrückt, öffnet sich ein Fenster, in welchem die Farben der Haare, der
Haut, des Hoodies und der Hose der Spielfigur angepasst werden kann. Diese Farben können durch
selbiges Fenster ebenfalls als Datei gespeichert bzw. wieder geladen werden.

\subsection{Minigames}

Die Minispiele werden alle durch das Drücken eines Start-Knopfes, welcher sich in dem zu der Fachschaft gehörenden Raum befindet.
Die dadurch erscheinende Level-Auswahl kann durch das in der oberen, rechten Ecke platzierte rote \glqq{}X\grqq{} verlassen werden.
In der Mitte der Szene finden sich drei Knöpfe mit den Aufschriften \glqq Level 1\grqq, \glqq Level 2\grqq und \glqq Level 3\grqq, welceh das entsprechende Level starten.
Jedes Level besitzt neben dem Knopf zum verlassen des Minispieles, ebenfalls einen Knopf, welcher drei gestapelte Balken abbildet und die Spielenden zurück zur Level-Auswahl führt.
Diese beiden Optionen bieten sich auch nach dem Abschließen eines Levels.
Zusätzlich gibt es hier noch die Möglichkeit des Neustarts, welche durch einen \glqq $\circlearrowright$\grqq{}-Knopf ermöglicht wird.

\subsubsection{Informatik}

Beim Minispiel der Fachschaft Informatik können die, sich auf der linken Seite des Bildschirms befindenden, Logikgatter mit der linken Maustaste angeklickt werden.
Bei Level 2 und Level 3 können mehrere ausgewählt werden.
Sollte ein bereits ausgewähltes erneut geklickt werden, wird die Auswahl zurückgenommen.
Ein neu ausgewähltes Gatter wird immer zu dem ersten freien Gatter, beginnend bei dem mit der kleinsten Nummer.
Falls also $G_1, G_2 \text{und} G_3$ ausgewählt und dann $G_1 \text{und} G_2$ wieder zurückgenommen werden, wird das danach gewählte Gatter zu $G_1$ statt $G_2$.

\subsubsection{Mathematik}

Beim Minispiel der Fachschaft Mathematik können die Tangramsteine mittels der linken Maustaste angeklickt werden.
Solange diese Taste gedrückt bleibt, folgt der Stein dem Mauszeiger und kann so verschoben werden.
Während dieses Zustandes kann die Taste \glqq R \grqq{} gedrückt werden, wodurch der Stein um $45\deg$ im Uhrzeigersinn gedreht wird.

\subsubsection{Chemie}

\subsubsection{Physik}

\subsubsection{Biologie}


\subsection{Kampf}


\section{Tests}

\subsection{\texttt{Unit Testing}}

Mithilfe des \texttt{JUnit}-Frameworks wurde eine Reihe von Tests implementiert,
welche die Funktionalität des Szenenstapels, also dessen Operationen, sowie das
Szenen-Verdeckungs-Einstellungs-System testet. Die Tests werden bei jedem Push
auf die GitHub-Repository ausgeführt, es ist also gleich ersichtlich, wenn eine
Änderung unbeabsichtigte Nebeneffekte hatte.

\subsection{Testen während der Entwicklung}

Bei der Implementierung von neuen Funktionen wurden diese stets umgehend soweit
getestet und ggf. geändert, dass sie allgemein dem Erwartungsbild entsprachen.
Andernfalls wurden einige Fehler im späteren Verlauf der Entwicklung beim Spielen
des Spiels aufgedeckt, gemeldet und, soweit möglich, behoben.
